
% ============================================================
% SECTION 02
% INVARIANCE AND IDENTITY
% Chapter 02: Invariance
% Book I: Ontology of Viability
% ============================================================

\section{Invariance and Identity}
\label{sec:invariance-and-identity}

% ------------------------------------------------------------
% 2.1. THE NEED TO DISTINGUISH IDENTITY FROM INVARIANCE
% ------------------------------------------------------------

The relationship between invariance and identity is
foundational to the conceptual development of this chapter.
Both concepts concern the determinations of an existent,
but they address different questions and must not be treated
as interchangeable.

The preceding chapter established that being and identity
are ontologically inseparable. To exist is to be
identity-bearing; identity is not an additional property
that is subsequently attached to an otherwise
identity-neutral existent.

The present chapter does not revise this commitment.
Instead, it examines how invariance can be understood
in relation to an existent whose identity is already
ontologically inseparable from its being.

Identity concerns the identity-bearing determination
of an existent as the existent that it is. Invariance,
according to the provisional relational definition
adopted here, concerns the relation between an existent
or one of its determinate aspects and a specified
domain of changes, in relation to which something
remains unchanged in a particular respect.

The distinction is therefore not between two
independently existing entities. It is between two
conceptual questions concerning the same existent.

The question of identity cannot be replaced by the
question of invariance, and the question of invariance
cannot be answered merely by repeating the proposition
that an existent is identical with itself.

% ------------------------------------------------------------
% 2.2. IDENTITY DOES NOT REQUIRE AN IMMEDIATE
% CLAIM OF INVARIANCE
% ------------------------------------------------------------

The inseparability of being and identity does not,
by itself, establish which determinations of an
existent remain invariant under a specified domain
of change.

To affirm that an existent is identity-bearing is
not yet to identify the aspects that remain unchanged
when that existent undergoes transformation.

For example, consider an existent whose spatial
position, external configuration, or relational
arrangement changes. The fact that these changes
occur does not automatically establish that its
identity has been lost. However, neither does the
continuation of its identity automatically identify
a particular feature that remains invariant.

These are distinct questions:

\begin{enumerate}
    \item Whether the existent is identity-bearing.
    \item Which of its determinations change under
    a specified transformation.
    \item Which of its determinations remain unchanged
    in relation to that transformation.
    \item Whether the identified invariance has any
    bearing on the persistence or continuity of its identity.
\end{enumerate}

The fourth question cannot be answered in advance.
It requires an argument concerning the relationship
between the particular invariant determination and
the identity of the existent.

Consequently, identity must not be reduced to a
collection of features that happen to remain unchanged
within a particular domain of observation or transformation.

% ------------------------------------------------------------
% 2.3. INVARIANCE DOES NOT CONSTITUTE IDENTITY
% ------------------------------------------------------------

A further distinction concerns the temptation to
treat invariance as the foundation or constitutive
source of identity.

If an existent preserves a particular feature under
a specified transformation, it may be tempting to
conclude that this feature is what makes the existent
the existent that it is.

Such a conclusion is not justified by the observation
of invariance alone.

The fact that a determination remains unchanged
under a particular domain of transformations establishes
a relational claim about that determination. It does
not establish that the determination is the ontological
ground of the existent's identity.

The distinction becomes especially important when
the invariant feature is structural, relational,
functional, or quantitative. The persistence of such
a feature may be relevant to the analysis of identity,
but its relevance does not automatically make it
the source of identity.

The ontology of viability therefore avoids identifying
an invariant determination with the identity of the
existent unless a separate argument establishes
that relationship.

The origin and constitution of identity remain
questions that require further ontological development.
They cannot be resolved merely through the relational
definition of invariance.

% ------------------------------------------------------------
% 2.4. THE SAME EXISTENT AND THE SAME DETERMINATION
% ------------------------------------------------------------

A distinction must also be maintained between the
identity of an existent and the persistence of a
particular determination belonging to it.

An existent may be considered under different
conditions or transformations. In such cases,
the question of whether the existent remains the
same must be distinguished from the question of
whether a particular determination remains unchanged.

These questions may have different answers.

A determination may remain invariant while other
determinations change. Conversely, a particular
determination may change while the existent's
identity remains an open question or continues
without that determination.

It is therefore necessary to distinguish at least
three analytical levels:

\begin{enumerate}
    \item The identity-bearing existence of the
    existent as such.

    \item The particular determinations through
    which the existent is considered.

    \item The invariance or variation of those
    determinations in relation to a specified
    domain of change.
\end{enumerate}

This distinction does not imply that the existent
can be ontologically separated from its determinations.
It only prevents the identity of the existent from
being equated with the unchanged persistence of
every particular determination.

The ontological relationship between the identity
of an existent and the determinations through which
it is constituted remains a subject for further
investigation.

% ------------------------------------------------------------
% 2.5. INVARIANCE AND THE FORMAL PRINCIPLE OF
% SELF-IDENTITY
% ------------------------------------------------------------

Formal self-identity is commonly expressed through
the proposition:

\[
A = A
\]

This proposition expresses the formal identity of
a term with itself. It does not, by itself, establish
that every determination associated with the referent
of that term remains unchanged under every possible
transformation.

The formal expression of self-identity and the
relational formulation of invariance therefore
operate at different conceptual levels.

Self-identity concerns the formal relation of a
term to itself. Invariance concerns whether a
specified determination remains unchanged under
a specified domain of transformations.

The distinction can be represented schematically.
Let \(A\) denote an existent and let \(T\) denote
a transformation within a specified domain.
If \(F\) represents a particular determination
of \(A\), an invariance claim may be expressed as:

\[
F(T(A)) = F(A)
\]

This expression is only a schematic representation.
It assumes that \(F\), \(T\), and the relevant
comparison are sufficiently defined. It does not
establish that the same mathematical formulation
applies to every ontological case.

The expression indicates that the determination
under consideration remains unchanged under the
specified transformation. It does not establish
that the entire existent is unchanged, nor that
\(F\) constitutes its identity.

Formal self-identity and invariance must therefore
remain conceptually distinct.

% ------------------------------------------------------------
% 2.6. THE POSSIBILITY OF MULTIPLE INVARIANT
% DETERMINATIONS
% ------------------------------------------------------------

An existent may be examined in relation to different
domains of change and different respects of
determination. Consequently, an invariance claim
concerning one aspect does not automatically
determine the status of every other aspect.

For example, a particular transformation may
preserve a structural relation while altering
a spatial configuration. Another transformation
may preserve a quantitative determination while
changing a relational arrangement.

These possibilities show that invariance claims
must be restricted to the specific determinations
and transformation domains under consideration.

The analysis does not yet establish whether
multiple invariant determinations are independent,
mutually dependent, hierarchically ordered, or
constitutively related.

Nor does the identification of multiple invariant
determinations establish that an existent possesses
multiple identities.

The question of how identity may relate to different
determinations across different ontological levels
remains unresolved. It must not be settled merely
by multiplying the number of invariant aspects
identified in an analysis.

This restriction is particularly important for
the subsequent development of the ontology of
viability, where the relationship between identity,
determination, and the different realms of being
will require a more comprehensive account.

% ------------------------------------------------------------
% 2.7. INVARIANCE AND THE LIMITS OF IDENTITY
% CLAIMS
% ------------------------------------------------------------

The distinction between identity and invariance
also imposes limits on what can be inferred from
the observation of an invariant feature.

First, the persistence of a feature under one
specified transformation does not establish its
persistence under every possible transformation.

Second, the persistence of a feature does not
establish that it is essential to the identity
of the existent.

Third, the alteration of a feature does not
establish that the existent has lost its identity.

Fourth, the absence of an identified invariant
feature does not establish that the existent
lacks identity or that no invariance is present.

These restrictions prevent the analysis from
confusing the limits of an invariance claim with
the ontological status of identity.

An identity-bearing existent cannot be treated
as ontologically dependent upon the observer's
ability to identify its invariant determinations.
The observer may recognize some aspects while
remaining unaware of others.

This epistemic limitation must be acknowledged
without converting it into an ontological
conclusion.

% ------------------------------------------------------------
% 2.8. THE RELATIONSHIP BETWEEN IDENTITY AND
% INVARIANCE REMAINS OPEN
% ------------------------------------------------------------

The analysis so far establishes a distinction
between identity and invariance, but it does not
establish that the two are unrelated.

An invariant determination may be relevant to
the way an existent is identified, distinguished,
or understood across different conditions.
It may also be relevant to the examination of
continuity and persistence.

However, the precise ontological significance
of such a determination cannot be inferred
from invariance alone.

Several questions remain open:

\begin{itemize}
    \item Can a particular invariant determination
    be constitutive of an existent's identity?

    \item Can an existent remain identity-bearing
    while all the determinations examined within
    a particular domain undergo change?

    \item Is there a necessary relationship between
    identity and any specific form of invariance?

    \item Can different forms of invariance have
    different ontological roles?

    \item How should the relationship between
    identity and invariance be understood across
    different ontological realms?
\end{itemize}

These questions require further argument. The
present chapter does not answer them by assuming
that invariance is either wholly external to
identity or necessarily constitutive of it.

The methodological commitment is to preserve
the distinction while keeping their possible
ontological relationship open.

% ------------------------------------------------------------
% 2.9. SYNTHESIS
% ------------------------------------------------------------

The relationship between invariance and identity
can now be stated in a preliminary form.

Identity concerns the identity-bearing existence
of an existent. Invariance concerns the relation
between that existent, or one of its determinate
aspects, and a specified domain of changes,
in relation to which something remains unchanged
in a particular respect.

The two concepts are therefore distinct but
potentially connected.

The distinction prevents three reductions:

\begin{enumerate}
    \item Identity must not be reduced to the
    persistence of invariant features.

    \item Invariance must not be treated as the
    ontological source of identity.

    \item The recognition of invariance must not
    be confused with the existence of identity.
\end{enumerate}

At the same time, the analysis does not deny
that invariant determinations may contribute
to a more comprehensive explanation of identity.
It only requires that such a contribution be
established through further argument.

The foundational commitment that being and
identity are inseparable remains intact. The
relational definition of invariance operates
within this commitment and does not replace it.

% ------------------------------------------------------------
% 2.10. TRANSITION
% ------------------------------------------------------------

Having distinguished invariance from identity,
the analysis must now examine the relationship
between invariance and difference.

The concept of invariance requires the specification
of a domain of change or transformation. Such
a domain necessarily raises questions about
difference: what varies, what remains unchanged,
and in what respect can these conditions be
distinguished?

The following section examines whether invariance
can be understood without a clear account of
difference, and whether the persistence of a
determination in relation to some differences
can coexist with variation in other respects.

% ============================================================
% END OF SECTION 02
% ============================================================
